% !TeX root = ../../../main.tex
\section{The Rigidified Picard Scheme and the Abel--Jacobi Map}\label{section:generalized-picard-and-abel-jacobi}

Our presentation of the rigidified Picard scheme and the Abel--Jacobi
morphism follows Guignard \cite[\S4]{Guignard2018} and
Takeuchi \cite[\S3]{takeuchi2018blow}. For background on the classical
theory of Jacobians and relative Picard schemes, we refer respectively to
Milne \cite{milneAV} and Bosch--L\"utkebohmert--Raynaud
\cite[Chapter~9]{BLR1990}. Precise references for the individual results
are given where they are used.

\subsection*{The rigidified Picard scheme}
Although the main theorem is over the fixed perfect field $k$, the
rigidified Picard construction is most naturally stated over an arbitrary
base scheme $S$.
Let $S$ be a scheme and let $f:C\to S$ be a projective smooth morphism whose geometric fibers are connected and of dimension 1. Let $\mdls$ be a modulus on $C$, defined as an effective Cartier divisor of $C/S$
(i.e., a closed subscheme of $C$ which is finite flat of finite presentation over $S$).
We denote the projection $C \times_S T \to T$ by $\operatorname{pr}$ for any $S$-scheme $T$.

Let $d$ be an integer. For an $S$-scheme $T$, we consider the set of data $(\cL, \psi)$ where:
\begin{itemize}
    \item $\cL$ is an invertible sheaf of degree $d$ on $C_T$.
    \item $\psi: \Oo_{\mdls_T} \xrightarrow{\sim} \cL|_{\mdls_T}$ is a trivialization of $\cL$ along the modulus.
\end{itemize}
Two such pairs $(\cL, \psi)$ and $(\cL', \psi')$ are said to be isomorphic if there exists an isomorphism of invertible sheaves $\alpha\colon \cL \to \cL'$ such that the following diagram commutes:
\[
\begin{tikzcd}
& \Oo_{\mdls_T} \arrow[dl, "\psi"'] \arrow[dr, "\psi'"] & \\
\cL|_{\mdls_T} \arrow[rr, "\alpha|_{\mdls_T}"'] & & \cL'|_{\mdls_T}
\end{tikzcd}
\]
We define the presheaf $\operatorname{Pic}_{C,\mdls}^{d,\mathrm{pre}}$ on
$\operatorname{Sch}/S$ by assigning to $T$ the set of isomorphism classes
of such pairs. Let $\PicCm[d]$ denote the \'etale sheafification of this
presheaf.

We use the following structural facts, recorded in
\cite[\S3, fundamental properties (1)--(4)]{takeuchi2018blow}. When $\mdls$ is everywhere
positive, representability and smoothness of the full rigidified Picard
presheaf are proved in \cite[Proposition~4.8]{Guignard2018}; when
$\mdls=0$, the corresponding results for the ordinary relative Picard
scheme are given in \cite[\S9.3, Theorem~1]{BLR1990}.
\begin{enumerate}
    \item $\PicCm[d]$ is represented by an $S$-scheme. If $\mdls$ is
    faithfully flat over $S$, the presheaf
    $\operatorname{Pic}_{C,\mdls}^{d,\mathrm{pre}}$ is already an \'etale
    sheaf.
    \item $\PicCm[0]$ is a smooth commutative group $S$-scheme with
    geometrically connected fibers.
    \item For any $d$, $\PicCm[d]$ is a $\PicCm[0]$-torsor.
\end{enumerate}
We call $\PicCm[0]$ the \emph{generalized Jacobian} of $C$ with modulus
$\mdls$. When $\mdls=0$, it is the ordinary Jacobian, and
$\PicCm[d]=\PicC[d]$.

The rigidified Picard schemes are naturally related as the modulus varies.
Indeed, let $\mdls_1\subset\mdls_2$. Restriction of rigidifications defines
a natural morphism of \'etale sheaves
\[
\rho_{\mdls_2,\mdls_1}\colon
\operatorname{Pic}_{C,\mdls_2}^d
\longrightarrow
\operatorname{Pic}_{C,\mdls_1}^d.
\]
When $\mdls_2$ is finite over $S$, this morphism is an epimorphism
\cite[\S3]{takeuchi2018blow}.
Taking $\mdls_1=0$ gives, for every
modulus $\mdls$, a natural epimorphism of \'etale sheaves
\[
\rho_{\mdls,0}\colon \PicCm[d]\longrightarrow\PicC[d].
\]

\subsection*{The Abel--Jacobi Morphism and Its Fibers}\label{section:AbelJacobi}


Let $U = C \setminus \mdls$ be the complement of the modulus in $C$.
The effective Cartier divisors of degree $d$ which are prime to $\mdls$ are parameterized by the symmetric power $\Sym_S^d (U) = U^{(d)}$ over $S$ (see \cite[Proposition~4.12]{Guignard2018} and \cite[Chapter~III, Theorem~3.13]{milneAV}).
For any such divisor $D \in U^{(d)}$, the associated line bundle $\Oo_C(D)$ admits a canonical trivialization along $\mdls$.
Specifically, the canonical section $1_D$ is regular and non-vanishing on $\mdls$ because $\operatorname{supp}(D) \cap \operatorname{supp}(\mdls) = \emptyset$.
Its restriction is a nowhere-vanishing section on $\mdls$, thereby
determining a trivialization
$\psi_D:\Oo_{\mdls}\xrightarrow{\sim}\Oo_C(D)|_{\mdls}$.
This is functorial in families and yields a morphism from the symmetric power
to the rigidified Picard scheme over $S$:
\begin{equation}
    \Phi_d\colon U^{(d)} \to \PicCm[d], \quad D \mapsto [(\Oo_C(D), \psi_D)],
\end{equation}

The results of Guignard and Takeuchi used below require the degree to be
sufficiently large.\footnote{The cohomological reason for this bound is that,
for a degree-$d$ line bundle $\cL$, it ensures that both $\cL$ and
$\cL(-\mdls)$ have fiberwise degree at least $2g-1$. Their first cohomology
then vanishes by Serre duality and a degree argument; Riemann--Roch together
with cohomology and base change implies that the corresponding pushforwards
are locally free of the expected ranks and commute with base change. See
\cite[Proposition~4.3]{Guignard2018}. The additional inequality
$d\geq\deg\mdls$ ensures that $d-\deg\mdls$ is nonnegative, which is relevant
when $g=0$.} Let $g$ denote the genus of $C$, viewed as a locally
constant function on $S$. From this point on, we impose the standing
condition
\begin{equation} \label{equation:degree-condition}
    d \geq \max\{2g - 1 + \deg \mdls, \deg \mdls\}.
\end{equation}
This is Takeuchi's condition \cite[\S3, (3.1)]{takeuchi2018blow} and is, in
particular, strong enough for the application of
\cite[Proposition~4.14]{Guignard2018} below. If $S$ is quasi-compact, the
locally constant functions $g$ and $\deg\mdls$ take only finitely many
values; hence a single integer $d$ satisfying this condition exists.

Suppose that $d$ satisfies \Cref{equation:degree-condition}. When $\mdls=0$
and $C\to S$ admits a section, the Abel--Jacobi morphism
\[
    C^{(d)}\longrightarrow\PicC[d]
\]
is a projective space bundle. In particular, it is proper and smooth with
geometrically connected fibers.

When $\mdls>0$, Guignard \cite[Proposition~4.14]{Guignard2018} proves the
stronger statement that
\[
    \Phi_d\colon U^{(d)}\longrightarrow\PicCm[d]
\]
is, for the fppf topology, a torsor under the additive group scheme of a
vector bundle of rank $d-\deg\mdls-g+1$ over $\PicCm[d]$. In particular,
$\Phi_d$ is fppf-locally isomorphic over $\PicCm[d]$ to an affine space of
relative dimension $d-\deg\mdls-g+1$; it is smooth and surjective, with
geometrically connected fibers.

We now return to the standing situation $S=\Spec(k)$, which remains in
force for the rest of the chapter. From Guignard's result, we retain only
the following fiberwise consequence, stated together with the classical
unramified case.
\begin{theorem}\label{theorem:abel-jacobi-fibers}
Assume that $S=\Spec(k)$ and that $d$ satisfies
\Cref{equation:degree-condition}. Let
$\bar s=\Spec\Omega\to\PicCm[d]$ be a geometric point, where $\Omega$ is
separably closed. Then the geometric fiber of the Abel--Jacobi morphism
\[
    \Phi_d\colon U^{(d)}\longrightarrow\PicCm[d]
\]
at $\bar s$ is isomorphic to
\[
\begin{cases}
\mathbb{A}_{\Omega}^{d - \deg \mdls - g + 1} & \text{if } \mdls > 0, \\
\mathbb{P}_{\Omega}^{d - g} & \text{if } \mdls = 0.
\end{cases}
\]
Thus every geometric fiber is an affine space when $\mdls>0$ and a
projective space when $\mdls=0$.
\end{theorem}
\begin{proof}
For $\mdls>0$, this follows from
\cite[Proposition~4.14]{Guignard2018}. For $\mdls=0$, see
\cite[Proposition~2.1.4]{Toth2011}.
Both cases are also stated in
\cite[Propositions~3.13--3.14]{tendler2015geometricclassfieldtheory}.
\end{proof}

